\documentclass[letterpaper,11pt]{article}

% --- PAQUETES DE IDIOMA ---
\usepackage[spanish,mexico,es-noshorthands]{babel}
\usepackage[utf8]{inputenc}
\usepackage{csquotes}
\usepackage{iftex}

% --- PAQUETES MATEMÁTICOS ---
\usepackage{amsmath}
\usepackage{amssymb}

% --- PAQUETES GRÁFICOS Y DE DISEÑO ---
\usepackage{graphicx}
\usepackage{float}
\usepackage{placeins}
\usepackage[table]{xcolor}
\usepackage{makecell}
\usepackage{multirow}
\usepackage{array}

% --- LISTAS ---
\usepackage{enumitem}

% --- PAQUETES DE TABLAS ---
\usepackage{booktabs}
\usepackage{longtable}

% --- COLORES POR FUNCIÓN DEL NIST CSF 2.0 ---
\definecolor{csfGV}{HTML}{FFF2CC}       % Govern - Amarillo suave
\definecolor{csfID}{HTML}{D9E1F2}       % Identify - Azul suave
\definecolor{csfPR}{HTML}{E2EFDA}       % Protect - Verde suave
\definecolor{csfDE}{HTML}{FCE4D6}       % Detect - Naranja suave
\definecolor{csfRS}{HTML}{F8CBAD}       % Respond - Coral suave
\definecolor{csfRC}{HTML}{D0CECE}       % Recover - Gris suave
\definecolor{xlSummary}{HTML}{CFE2F3}   % Azul claro para encabezados

% --- RUTAS DE IMÁGENES ---
\graphicspath{{portada_img/}{../portada_img/}}

% --- ÍNDICE CON PUNTOS ---
\usepackage{tocloft}
\renewcommand{\cftsecleader}{\cftdotfill{\cftdotsep}}

% --- CONFIGURACIÓN DE PÁGINA (GEOMETRÍA) ---
\usepackage[
    left=18mm, 
    right=18mm,
    top=30mm,
    bottom=24mm,
    headheight=77pt, 
    papersize={21.59cm,27.94cm}
]{geometry}

% --- FORMATO DE TEXTO ---
\usepackage{setspace}
\onehalfspacing
\usepackage{parskip}

% --- ENCABEZADOS Y PIES DE PÁGINA ---
\usepackage{fancyhdr}
\pagestyle{fancy}
\fancyhf{} 
\fancyhead[L]{\includegraphics[height=1.5cm]{portada_img/der.png}}
\fancyhead[R]{\includegraphics[height=1.5cm]{portada_img/izq.png}}
\fancyfoot[C]{\thepage} 
\renewcommand{\headrulewidth}{0.4pt} 
\renewcommand{\footrulewidth}{0pt}  

% --- BIBLIOGRAFÍA ---
\usepackage[backend=biber, style=apa, sortcites, url=true]{biblatex}
\addbibresource{referencias.bib}

% --- HIPERVÍNCULOS ---
\usepackage{hyperref}
\hypersetup{
    colorlinks=true,
    linkcolor=blue,
    citecolor=blue,
    filecolor=magenta,
    urlcolor=cyan
}

% Tipos de columna personalizados
\newcolumntype{L}[1]{>{\raggedright\arraybackslash}p{#1}}
\newcolumntype{C}[1]{>{\centering\arraybackslash}p{#1}}
\newcolumntype{R}[1]{>{\raggedleft\arraybackslash}p{#1}}
\newcolumntype{M}[1]{>{\centering\arraybackslash}m{#1}}

\begin{document}

\thispagestyle{empty} % Sin encabezado ni pie en la portada

% --- ENCABEZADO DE PORTADA (TABLA DE 3 COLUMNAS) ---
\begin{center}
    % La tabla usa 'm' para centrar verticalmente las imágenes y el texto
    \begin{tabular}{m{0.15\textwidth} m{0.60\textwidth} m{0.15\textwidth}}
        % Logo Izquierdo (UNAM)
        \includegraphics[width=2cm]{portada_img/der.png} & 
        
        % Título Central
        \centering\Huge\textbf{Universidad Nacional Autónoma de México} &
        
        % Logo Derecho (Facultad)
        \raggedleft\includegraphics[width=2cm]{portada_img/izq.png}
    \end{tabular}
\end{center}

\vspace{0.3cm}

% --- INFORMACIÓN DEL REPORTE ---
\begin{center}
    {\LARGE \textbf{Facultad de Ingeniería}}\\[0.8cm]

    {\LARGE \textbf{Alumno:}}\\[0.3cm]
    {\Large Lara Hernandez Angel Husiel - 320060829} \\[0.6cm]
    {\Large Lugo Manzano Rodrigo - 320206968} \\[0.6cm]
    {\Large Flores Colin Victor Jaziel - 320266083} \\[0.6cm]
    {\Large Garcia Cortes Adolfo de Jesus - 320317264} \\[0.8cm]

    {\LARGE \textbf{Seguridad informatica basica}} \\[0.8cm]

    {\LARGE Semestre: 2027-1} \\[0.8cm]
    
    {\LARGE \textbf{Actividad 8}}\\[0.4cm]

    {\LARGE {Mapeo de Categorías y Subcategorías de The NIST Cybersecurity Framework (CSF) 2.0 con las Actividades del Semestre}}\\[0.8cm]

    {\Large \textbf{Fecha de Entrega:}}\\[0.2cm]
    {\Large 14 de Noviembre del 2026}
\end{center}


\newpage
{\hypersetup{linkcolor=black}\tableofcontents}

\newpage
\setcounter{page}{1}
\singlespacing

% =========================================================
\section{Categorías y Subcategorías por Función Core de The NIST CSF 2.0}
% =========================================================

El núcleo (\textit{Core}) de \textbf{The NIST Cybersecurity Framework (CSF) 2.0} \autocite{nist-csf-20} organiza la gestión de riesgos de ciberseguridad en \textbf{6 Funciones Core}, las cuales se dividen en \textbf{22 Categorías} y \textbf{106 Subcategorías}. En la Tabla~\ref{tab:estructura-csf} se presentan las categorías y subcategorías consideradas para cada función del framework.

\begin{table}[H]
\centering
\footnotesize
\setlength{\tabcolsep}{5pt}
\renewcommand{\arraystretch}{1.24}
\caption{Funciones Core, Categorías y Subcategorías consideradas en The NIST CSF 2.0.}
\label{tab:estructura-csf}
\begin{tabular}{|L{3.0cm}|L{7.2cm}|C{5.8cm}|}
\hline
\rowcolor{blue!20}
\textbf{Función Core} & \textbf{Categorías Consideradas (22)} & \textbf{Subcategorías Consideradas (106)} \\
\hline

\textbf{1. GOVERN (\texttt{GV})}\newline (Gobernanza) &
\textbf{GV.OC}: Contexto Organizacional\newline
\textbf{GV.RM}: Estrategia de Gestión de Riesgos\newline
\textbf{GV.RR}: Roles, Responsabilidades y Autoridades\newline
\textbf{GV.PO}: Políticas\newline
\textbf{GV.OV}: Supervisión (\textit{Oversight})\newline
\textbf{GV.SC}: Gestión de Riesgos en Cadena de Suministro &
\texttt{GV.OC-01} a \texttt{GV.OC-05} (5 subcategorías)\newline
\texttt{GV.RM-01} a \texttt{GV.RM-07} (7 subcategorías)\newline
\texttt{GV.RR-01} a \texttt{GV.RR-04} (4 subcategorías)\newline
\texttt{GV.PO-01} a \texttt{GV.PO-02} (2 subcategorías)\newline
\texttt{GV.OV-01} a \texttt{GV.OV-03} (3 subcategorías)\newline
\texttt{GV.SC-01} a \texttt{GV.SC-10} (10 subcategorías) \\
\hline

\textbf{2. IDENTIFY (\texttt{ID})}\newline (Identificar) &
\textbf{ID.AM}: Gestión de Activos (\textit{Asset Management})\newline
\textbf{ID.RA}: Evaluación de Riesgos (\textit{Risk Assessment})\newline
\textbf{ID.IM}: Mejora (\textit{Improvement}) &
\texttt{ID.AM-01} a \texttt{ID.AM-08} (7 subcategorías)\newline
\texttt{ID.RA-01} a \texttt{ID.RA-10} (10 subcategorías)\newline
\texttt{ID.IM-01} a \texttt{ID.IM-04} (4 subcategorías) \\
\hline

\textbf{3. PROTECT (\texttt{PR})}\newline (Proteger) &
\textbf{PR.AA}: Gestión de Identidad, Autenticación y Acceso\newline
\textbf{PR.AT}: Concientización y Capacitación\newline
\textbf{PR.DS}: Seguridad de los Datos (\textit{Data Security})\newline
\textbf{PR.PS}: Seguridad de Plataformas\newline
\textbf{PR.IR}: Resiliencia de la Infraestructura Tecnológica &
\texttt{PR.AA-01} a \texttt{PR.AA-06} (6 subcategorías)\newline
\texttt{PR.AT-01} a \texttt{PR.AT-02} (2 subcategorías)\newline
\texttt{PR.DS-01}, \texttt{02}, \texttt{10}, \texttt{11} (4 subcategorías)\newline
\texttt{PR.PS-01} a \texttt{PR.PS-06} (6 subcategorías)\newline
\texttt{PR.IR-01} a \texttt{PR.IR-04} (4 subcategorías) \\
\hline

\textbf{4. DETECT (\texttt{DE})}\newline (Detectar) &
\textbf{DE.CM}: Monitoreo Continuo\newline
\textbf{DE.AE}: Análisis de Eventos Adversos &
\texttt{DE.CM-01}, \texttt{02}, \texttt{03}, \texttt{06}, \texttt{09} (5 subcategorías)\newline
\texttt{DE.AE-02} a \texttt{DE.AE-08} (6 subcategorías) \\
\hline

\textbf{5. RESPOND (\texttt{RS})}\newline (Responder) &
\textbf{RS.MA}: Gestión de Incidentes\newline
\textbf{RS.AN}: Análisis de Incidentes\newline
\textbf{RS.CO}: Reporte y Comunicación de Respuesta\newline
\textbf{RS.MI}: Mitigación de Incidentes &
\texttt{RS.MA-01} a \texttt{RS.MA-05} (5 subcategorías)\newline
\texttt{RS.AN-03}, \texttt{06}, \texttt{07}, \texttt{08} (4 subcategorías)\newline
\texttt{RS.CO-02} a \texttt{RS.CO-03} (2 subcategorías)\newline
\texttt{RS.MI-01} a \texttt{RS.MI-02} (2 subcategorías) \\
\hline

\textbf{6. RECOVER (\texttt{RC})}\newline (Recuperar) &
\textbf{RC.RP}: Ejecución del Plan de Recuperación\newline
\textbf{RC.CO}: Comunicación de Recuperación &
\texttt{RC.RP-01} a \texttt{RC.RP-06} (6 subcategorías)\newline
\texttt{RC.CO-03} a \texttt{RC.CO-04} (2 subcategorías) \\
\hline
\end{tabular}
\end{table}

% =========================================================
\section{Mapeo de Categorías, Subcategorías y Actividades Realizadas en el Semestre}
% =========================================================

A continuación se presenta el mapeo de las categorías y subcategorías del \textbf{NIST CSF 2.0} que se relacionan directamente con los temas explicados en clase, las \textbf{Actividades 1 a 7}, las \textbf{Prácticas 1 a 7}, la \textbf{Presentación en clase} y las \textbf{Tareas 1 a 3} realizadas durante el semestre. Las subcategorías que pertenecen a una misma categoría están unificadas bajo una sola celda en la primera columna.

% ---------------------------------------------------------
\subsection{Función 1: GOVERN (GV) --- Gobernanza (11 Subcategorías)}
% ---------------------------------------------------------

\begingroup
\scriptsize
\setlength{\tabcolsep}{4pt}
\renewcommand{\arraystretch}{1.24}
\begin{longtable}{|L{2.8cm}|L{4.5cm}|L{2.9cm}|L{6.5cm}|}
\caption{Mapeo de Categorías y Subcategorías de la función \textbf{GOVERN (GV)} con las actividades del semestre.}
\label{tab:mapeo-gv} \\
\hline
\rowcolor{csfGV}
\textbf{Categoría} & \textbf{Subcategoría Relacionada} & \textbf{Actividad / Práctica / Tarea} & \textbf{Qué se hizo en el semestre y cómo se relaciona} \\
\hline
\endfirsthead
\hline
\rowcolor{csfGV}
\textbf{Categoría} & \textbf{Subcategoría Relacionada} & \textbf{Actividad / Práctica / Tarea} & \textbf{Qué se hizo en el semestre y cómo se relaciona} \\
\hline
\endhead
\hline
\endfoot

% --- GV.OC (4 subcategorías) ---
\textbf{GV.OC}\newline Contexto Organizacional &
\textbf{\texttt{GV.OC-01}}: La misión organizacional se comprende e informa la gestión de riesgos. &
\textbf{Clase y Actividad 2} &
Se analizó el giro de negocio, misión, visión y objetivos estratégicos de una organización para comprender qué procesos operativos deben protegerse prioritariamente. \\
\cline{2-4}\nopagebreak
 &
\textbf{\texttt{GV.OC-02}}: Las partes interesadas internas y externas y sus expectativas de ciberseguridad se comprenden. &
\textbf{Actividad 2, Actividad 6 y Tarea 2} &
Se identificaron las áreas internas y externas (clientes, socios comerciales, reguladores), así como las expectativas éticas de la sociedad en la \textbf{Tarea 2}. \\
\cline{2-4}\nopagebreak
 &
\textbf{\texttt{GV.OC-03}}: Los requisitos legales, regulatorios y contractuales se comprenden y gestionan. &
\textbf{Clase, Actividad 1 y Actividad 7} &
En la \textbf{Actividad 1} se revisaron normativas como ISO/IEC 27001, y en la \textbf{Actividad 7} se evaluó el impacto regulatorio y contractual (protección de datos personales, ISO/SAE 21434, TISAX y acuerdos NDA). \\
\cline{2-4}\nopagebreak
 &
\textbf{\texttt{GV.OC-05}}: Los servicios y capacidades de los que depende la organización se comprenden. &
\textbf{Práctica 2 y Actividad 3} &
En la \textbf{Práctica 2}, mediante consultas DNS y \texttt{whois}, se identificó que dominios institucionales dependen de servicios alojados en Amazon Web Services (AWS) y Cloudflare. \\
\hline

% --- GV.RM (4 subcategorías) ---
\textbf{GV.RM}\newline Estrategia de Gestión de Riesgos &
\textbf{\texttt{GV.RM-02}}: Las declaraciones de apetito y tolerancia al riesgo se establecen y mantienen. &
\textbf{Actividad 7 y Tarea 3} &
En la \textbf{Actividad 7} se definieron umbrales de tolerancia en indicadores KRI, y en la \textbf{Tarea 3} (conferencia \textit{The Defensible Minimum}) se estudió la aceptación consciente del riesgo residual. \\
\cline{2-4}\nopagebreak
 &
\textbf{\texttt{GV.RM-03}}: La gestión de riesgos de ciberseguridad se incluye en los procesos de riesgo empresarial (ERM). &
\textbf{Clase y Actividad 7} &
En la \textbf{Actividad 7} se clasificó el riesgo dentro de categorías empresariales (Estratégico, Operativo, Financiero, Cumplimiento y Tecnológico) evaluando pérdidas económicas, costos de respuesta, reputación y multas. \\
\cline{2-4}\nopagebreak
 &
\textbf{\texttt{GV.RM-04}}: Se establece dirección estratégica que describe opciones de respuesta al riesgo. &
\textbf{Clase, Actividad 6 y Actividad 7} &
Se aplicaron en las \textbf{Actividades 6 y 7} las cuatro estrategias de respuesta al riesgo (\textbf{Mitigar, Transferir, Evitar y Aceptar}), justificando cuál elegir en cada escenario. \\
\cline{2-4}\nopagebreak
 &
\textbf{\texttt{GV.RM-06}}: Se establece un método estandarizado para calcular, documentar, categorizar y priorizar los riesgos. &
\textbf{Actividad 4, Actividad 6 y Actividad 7} &
En la \textbf{Actividad 4} se usó \textbf{CVSS v3.1}; en la \textbf{Actividad 6} se construyó la plantilla del \textbf{Registro de Riesgos} (Forma 1 y Forma 2); y en la \textbf{Actividad 7} se aplicó la fórmula cuantitativa de Nivel de Riesgo ($R = F \times I$). \\
\hline

% --- GV.RR (2 subcategorías) ---
\textbf{GV.RR}\newline Roles, Responsabilidades y Autoridades &
\textbf{\texttt{GV.RR-01}}: Liderazgo responsable y cultura ética consciente del riesgo. &
\textbf{Tarea 2} &
En la \textbf{Tarea 2} se analizó la responsabilidad y ética profesional en ciberseguridad frente a los mercados de exploits y la divulgación responsable de vulnerabilidades. \\
\cline{2-4}\nopagebreak
 &
\textbf{\texttt{GV.RR-02}}: Roles y responsabilidades de gestión de riesgos establecidos. &
\textbf{Actividad 6 y Actividad 7} &
En las \textbf{Actividades 6 y 7} se asignó formalmente un \textbf{Propietario del Riesgo} (\textit{Risk Owner}) y un área responsable para cada escenario y acción del plan de mitigación. \\
\hline

% --- GV.OV (1 subcategoría) ---
\textbf{GV.OV}\newline Supervisión (\textit{Oversight}) &
\textbf{\texttt{GV.OV-03}}: El desempeño de la gestión de riesgos se evalúa y revisa mediante métricas (KRIs). &
\textbf{Actividad 7} &
En la \textbf{Actividad 7} se diseñaron \textbf{3 Indicadores Clave de Riesgo (KRIs)} con métricas cuantitativas, frecuencia de medición y umbrales de alerta para evaluar los controles propuestos. \\
\hline
\end{longtable}
\endgroup

% ---------------------------------------------------------
\subsection{Función 2: IDENTIFY (ID) --- Identificar (12 Subcategorías)}
% ---------------------------------------------------------

\begingroup
\scriptsize
\setlength{\tabcolsep}{4pt}
\renewcommand{\arraystretch}{1.24}
\begin{longtable}{|L{2.8cm}|L{4.5cm}|L{2.9cm}|L{6.5cm}|}
\caption{Mapeo de Categorías y Subcategorías de la función \textbf{IDENTIFY (ID)} con las actividades del semestre.}
\label{tab:mapeo-id} \\
\hline
\rowcolor{csfID}
\textbf{Categoría} & \textbf{Subcategoría Relacionada} & \textbf{Actividad / Práctica / Tarea} & \textbf{Qué se hizo en el semestre y cómo se relaciona} \\
\hline
\endfirsthead
\hline
\rowcolor{csfID}
\textbf{Categoría} & \textbf{Subcategoría Relacionada} & \textbf{Actividad / Práctica / Tarea} & \textbf{Qué se hizo en el semestre y cómo se relaciona} \\
\hline
\endhead
\hline
\endfoot

% --- ID.AM (5 subcategorías) ---
\textbf{ID.AM}\newline Gestión de Activos &
\textbf{\texttt{ID.AM-01}}: Se mantienen inventarios de hardware gestionados por la organización. &
\textbf{Práctica 3, Práctica 6 y Actividad 3} &
En la \textbf{Práctica 3} se descubrieron hosts activos con \texttt{nmap -sP}, en la \textbf{Práctica 6} se creó un script de barrido para descubrir servidores en la red interna, y en la \textbf{Actividad 3} se elaboró el inventario de hardware y endpoints. \\
\cline{2-4}\nopagebreak
 &
\textbf{\texttt{ID.AM-02}}: Se mantienen inventarios de software, servicios y sistemas. &
\textbf{Práctica 2, Práctica 3, Práctica 7 y Actividad 3} &
En la \textbf{Práctica 2} se enumeraron subdominios con \texttt{subfinder} y tecnologías web con \texttt{httpx}; en la \textbf{Práctica 3} se identificaron puertos y versiones con \texttt{nmap -sV -A}; en la \textbf{Práctica 7} se inventarió el sistema con \texttt{systeminfo}; y en la \textbf{Actividad 3} se documentó el catálogo de software y APIs. \\
\cline{2-4}\nopagebreak
 &
\textbf{\texttt{ID.AM-03}}: Se mantienen representaciones de la comunicación de red y flujos de datos. &
\textbf{Práctica 6, Práctica 7 y Actividad 3} &
En las \textbf{Prácticas 6 y 7} se inspeccionaron interfaces (\texttt{ifconfig -a}) y tablas de ruteo (\texttt{route -n}), y en la \textbf{Actividad 3} se diseñó el diagrama de arquitectura de red dividido en zonas de seguridad. \\
\cline{2-4}\nopagebreak
 &
\textbf{\texttt{ID.AM-07}}: Se mantienen inventarios de datos para tipos de datos designados. &
\textbf{Práctica 6, Práctica 7 y Actividad 3} &
En las \textbf{Prácticas 6 y 7} se localizaron archivos sensibles, respaldos y código fuente (\texttt{find}, \texttt{grep -rnI}), y en la \textbf{Actividad 3} se clasificaron las bases de datos e información crítica. \\
\cline{2-4}\nopagebreak
 &
\textbf{\texttt{ID.AM-08}}: Gestión del ciclo de vida de sistemas (detección de software EOL y \textit{Shadow AI}). &
\textbf{Práctica 2, Práctica 4, Actividad 6 y Actividad 7} &
En las \textbf{Prácticas 2 y 4} se detectaron sistemas fuera de soporte (PHP 5.x y Debian 8.11 Jessie EOL). En las \textbf{Actividades 6 y 7} se analizó el uso no autorizado de servicios de IA (\textbf{ShadowAI}). \\
\hline

% --- ID.RA (6 subcategorías) ---
\textbf{ID.RA}\newline Evaluación de Riesgos &
\textbf{\texttt{ID.RA-01}}: Las vulnerabilidades en los activos se identifican, validan y registran. &
\textbf{Práctica 2, Práctica 4 y Actividad 4} &
En la \textbf{Práctica 2} se escaneó un portal web con \texttt{nuclei}; en la \textbf{Práctica 4} se ejecutó \texttt{LinPEAS} detectando 13 CVEs de escalación en Linux; y en la \textbf{Actividad 4} se registraron 27 vulnerabilidades CVE con su código CWE y puntaje CVSS v3.1. \\
\cline{2-4}\nopagebreak
 &
\textbf{\texttt{ID.RA-02}}: Se recibe inteligencia de amenazas cibernéticas (CTI) de fuentes externas. &
\textbf{Práctica 2, Actividad 5 y Tarea 3} &
En la \textbf{Práctica 2} se integraron las APIs de VirusTotal y GitHub en \texttt{subfinder}, y en la \textbf{Tarea 3} se analizaron las novedades de \textbf{MITRE ATT\&CK v21} y \textbf{MITRE ATLAS}. \\
\cline{2-4}\nopagebreak
 &
\textbf{\texttt{ID.RA-03}}: Las amenazas internas y externas se identifican y registran. &
\textbf{Actividad 5 y Presentación en Clase} &
En la \textbf{Actividad 5} y en la \textbf{Presentación en clase} se investigaron los grupos de amenaza más relevantes y se diseñó una taxonomía propia de adversarios externos e internos. \\
\cline{2-4}\nopagebreak
 &
\textbf{\texttt{ID.RA-04}}: Se identifican y registran impactos potenciales y probabilidades. &
\textbf{Actividad 6 y Actividad 7} &
En la \textbf{Actividad 6} se redactaron 10 escenarios de riesgo y en la \textbf{Actividad 7} se estimó cuantitativamente la frecuencia anual y el impacto para calcular el nivel de riesgo. \\
\cline{2-4}\nopagebreak
 &
\textbf{\texttt{ID.RA-06}}: Las respuestas al riesgo se eligen, priorizan, planifican y rastrean (POA\&M). &
\textbf{Actividad 6 y Actividad 7} &
Se documentó el estado de mitigación en el Registro de Riesgos (\textbf{Actividad 6}) y se construyó el \textbf{Plan de Acción e Hitos (POA\&M)} con 5 controles y fechas compromiso (\textbf{Actividad 7}). \\
\cline{2-4}\nopagebreak
 &
\textbf{\texttt{ID.RA-08}}: Procesos para recibir, analizar y responder a divulgaciones de vulnerabilidades. &
\textbf{Tarea 2} &
En la \textbf{Tarea 2} se investigó el ciclo de las vulnerabilidades \textit{Zero-Day}, los mercados blanco, gris y negro, y los programas de divulgación responsable (\textit{Bug Bounty}). \\
\hline

% --- ID.IM (1 subcategoría) ---
\textbf{ID.IM}\newline Mejora (\textit{Improvement}) &
\textbf{\texttt{ID.IM-02}}: Se identifican mejoras a partir de pruebas de seguridad y \textit{Penetration Testing}. &
\textbf{Prácticas 2 a 7 y Tarea 3} &
En las \textbf{Prácticas 2 a 7} se realizó un ejercicio completo de \textit{Penetration Testing} en laboratorio (reconocimiento, escaneo, fuerza bruta, explotación local, cracking de hashes y pivoteo), identificando las debilidades a corregir. \\
\hline
\end{longtable}
\endgroup

% ---------------------------------------------------------
\subsection{Función 3: PROTECT (PR) --- Proteger (11 Subcategorías)}
% ---------------------------------------------------------

\begingroup
\scriptsize
\setlength{\tabcolsep}{4pt}
\renewcommand{\arraystretch}{1.24}
\begin{longtable}{|L{2.8cm}|L{4.5cm}|L{2.9cm}|L{6.5cm}|}
\caption{Mapeo de Categorías y Subcategorías de la función \textbf{PROTECT (PR)} con las actividades del semestre.}
\label{tab:mapeo-pr} \\
\hline
\rowcolor{csfPR}
\textbf{Categoría} & \textbf{Subcategoría Relacionada} & \textbf{Actividad / Práctica / Tarea} & \textbf{Qué se hizo en el semestre y cómo se relaciona} \\
\hline
\endfirsthead
\hline
\rowcolor{csfPR}
\textbf{Categoría} & \textbf{Subcategoría Relacionada} & \textbf{Actividad / Práctica / Tarea} & \textbf{Qué se hizo en el semestre y cómo se relaciona} \\
\hline
\endhead
\hline
\endfoot

% --- PR.AA (3 subcategorías) ---
\textbf{PR.AA}\newline Gestión de Identidad, Autenticación y Control de Acceso &
\textbf{\texttt{PR.AA-01}}: Identidades y credenciales gestionadas para usuarios y servicios. &
\textbf{Práctica 1, Práctica 2 y Práctica 5} &
En la \textbf{Práctica 1} se crearon cuentas de usuario estándar en Linux; en la \textbf{Práctica 2} se gestionaron llaves de API; y en la \textbf{Práctica 5} se auditaron los archivos \texttt{/etc/passwd} y \texttt{/etc/shadow}. \\
\cline{2-4}\nopagebreak
 &
\textbf{\texttt{PR.AA-03}}: Usuarios y servicios son autenticados (políticas de fortaleza de contraseñas). &
\textbf{Práctica 3, Práctica 5 y Práctica 7} &
En la \textbf{Práctica 3} se auditó el servicio SSH mediante fuerza bruta por diccionario con \texttt{Hydra}, y en la \textbf{Práctica 5} se auditaron hashes \texttt{SHA512-crypt} con \texttt{John the Ripper}, comprobando el riesgo de usar contraseñas débiles de diccionario. \\
\cline{2-4}\nopagebreak
 &
\textbf{\texttt{PR.AA-05}}: Permisos de acceso definidos bajo el principio de mínimo privilegio. &
\textbf{Práctica 1 y Práctica 4} &
En la \textbf{Práctica 1} se configuró \texttt{/etc/sudoers} para otorgar privilegios mínimos. En la \textbf{Práctica 4} se auditaron con \texttt{LinPEAS} binarios SUID de root (\texttt{pkexec}) y se explotó \texttt{CVE-2021-4034} logrando escalar a \texttt{uid=0(root)}. \\
\hline

% --- PR.DS (3 subcategorías) ---
\textbf{PR.DS}\newline Seguridad de los Datos &
\textbf{\texttt{PR.DS-01}}: Se protege la confidencialidad e integridad de los datos en reposo (\textit{data-at-rest}). &
\textbf{Práctica 5, Práctica 6 y Práctica 7} &
En la \textbf{Práctica 5} se verificó la integridad de archivos con \texttt{sha256sum} y se analizó el hashing en \texttt{/etc/shadow}. En la \textbf{Práctica 6} se hallaron contraseñas en texto claro en respaldos, y en la \textbf{Práctica 7} se calculó el hash \texttt{SHA256} de \texttt{bot.py} con \texttt{certutil}. \\
\cline{2-4}\nopagebreak
 &
\textbf{\texttt{PR.DS-02}}: Se protege la confidencialidad e integridad de los datos en tránsito (\textit{data-in-transit}). &
\textbf{Práctica 2, Práctica 5, Práctica 7 y Actividad 7} &
En la \textbf{Práctica 2} se auditaron protocolos TLS con \texttt{nuclei}; en la \textbf{Práctica 5} se usó transferencia cifrada \texttt{scp}; en la \textbf{Práctica 7} se creó un túnel cifrado SSH (\texttt{ssh -L}); y en la \textbf{Actividad 7} se definieron controles DLP y CASB. \\
\cline{2-4}\nopagebreak
 &
\textbf{\texttt{PR.DS-11}}: Las copias de seguridad (respaldos) se crean y protegen de forma segura. &
\textbf{Práctica 6 y Tarea 3} &
En la \textbf{Práctica 6} se inspeccionaron respaldos comprimidos (\texttt{tar -ztvf}) comprobando el riesgo de no cifrar las copias de seguridad, y en la \textbf{Tarea 3} se analizó el borrado de respaldos por ransomware. \\
\hline

% --- PR.PS (4 subcategorías) ---
\textbf{PR.PS}\newline Seguridad de Plataformas &
\textbf{\texttt{PR.PS-01}}: Se establecen y aplican prácticas de gestión de configuración. &
\textbf{Práctica 1, Práctica 4 y Actividad 7} &
En la \textbf{Práctica 1} se configuró el entorno Debian 13, en la \textbf{Práctica 4} se auditaron configuraciones inseguras con \texttt{LinPEAS}, y en la \textbf{Actividad 7} se definieron políticas GPO. \\
\cline{2-4}\nopagebreak
 &
\textbf{\texttt{PR.PS-02}}: El software se mantiene, actualiza y se reemplazan versiones EOL. &
\textbf{Práctica 1, Práctica 2, Práctica 4 y Actividad 4} &
En la \textbf{Práctica 1} se ejecutó la actualización del sistema con \texttt{apt update \&\& apt upgrade}, y en las \textbf{Prácticas 2 y 4} se evidenció la vulnerabilidad de sistemas fuera de soporte (EOL). \\
\cline{2-4}\nopagebreak
 &
\textbf{\texttt{PR.PS-04}}: Se generan registros de log y se disponibilizan para monitoreo continuo. &
\textbf{Tarea 1 y Tarea 3} &
En la \textbf{Tarea 3} (conferencia \textit{Obfuscated Tale}) se estudió la generación de telemetría con PowerShell (Eventos 4103 y 4104) y Sysmon (Eventos 1, 7 y 17) hacia el SIEM. \\
\cline{2-4}\nopagebreak
 &
\textbf{\texttt{PR.PS-06}}: Prácticas de desarrollo seguro de software integradas en el ciclo de vida (SDLC). &
\textbf{Práctica 2, Práctica 6, Práctica 7 y Actividad 7} &
En la \textbf{Práctica 2} se detectó la falta de \textit{Subresource Integrity} (SRI) en scripts web, y en la \textbf{Práctica 6} se identificaron credenciales expuestas dentro de respaldos de código fuente. \\
\hline

% --- PR.IR (1 subcategoría) ---
\textbf{PR.IR}\newline Resiliencia de la Infraestructura Tecnológica &
\textbf{\texttt{PR.IR-01}}: Las redes y entornos están protegidos y segmentados lógicamente. &
\textbf{Práctica 6, Práctica 7 y Actividad 3} &
En las \textbf{Prácticas 6 y 7} se comprobó en laboratorio el aislamiento entre dos segmentos de red distintos y el uso de un servidor con doble interfaz para hacer pivoteo mediante túneles SSH. \\
\hline
\end{longtable}
\endgroup

% ---------------------------------------------------------
\subsection{Función 4: DETECT (DE) --- Detectar (6 Subcategorías)}
% ---------------------------------------------------------

\begingroup
\scriptsize
\setlength{\tabcolsep}{4pt}
\renewcommand{\arraystretch}{1.24}
\begin{longtable}{|L{2.8cm}|L{4.5cm}|L{2.9cm}|L{6.5cm}|}
\caption{Mapeo de Categorías y Subcategorías de la función \textbf{DETECT (DE)} con las actividades del semestre.}
\label{tab:mapeo-de} \\
\hline
\rowcolor{csfDE}
\textbf{Categoría} & \textbf{Subcategoría Relacionada} & \textbf{Actividad / Práctica / Tarea} & \textbf{Qué se hizo en el semestre y cómo se relaciona} \\
\hline
\endfirsthead
\hline
\rowcolor{csfDE}
\textbf{Categoría} & \textbf{Subcategoría Relacionada} & \textbf{Actividad / Práctica / Tarea} & \textbf{Qué se hizo en el semestre y cómo se relaciona} \\
\hline
\endhead
\hline
\endfoot

% --- DE.CM (3 subcategorías) ---
\textbf{DE.CM}\newline Monitoreo Continuo &
\textbf{\texttt{DE.CM-01}}: Las redes y los servicios de red se monitorean para encontrar eventos adversos. &
\textbf{Práctica 6, Práctica 7 y Actividad 7} &
En las \textbf{Prácticas 6 y 7} se verificaron conexiones y puertos activos, y en la \textbf{Actividad 7} se definió el monitoreo continuo de tráfico DNS/HTTPS hacia dominios no autorizados (KRI 1). \\
\cline{2-4}\nopagebreak
 &
\textbf{\texttt{DE.CM-03}}: La actividad del personal y la \textbf{tecnología de ciberengaño (Deception / Honeypots)} se monitorean continuamente. &
\textbf{Tarea 1 y Actividad 7} &
En la \textbf{Tarea 1} se estudió a fondo el despliegue y monitoreo continuo de \textbf{Honeypots}, \textit{Honey files}, \textit{Honey tokens} y cuentas señuelo en Active Directory para detectar intrusos y ataques de \textit{Kerberoasting}. \\
\cline{2-4}\nopagebreak
 &
\textbf{\texttt{DE.CM-09}}: El hardware, software, entornos de ejecución (\textit{runtime}) y datos se monitorean. &
\textbf{Práctica 6, Actividad 7 y Tarea 3} &
En la \textbf{Práctica 6} se inspeccionaron los procesos activos (\texttt{ps aux awf}), y en la \textbf{Tarea 3} se analizó el monitoreo en tiempo de ejecución con EDR y Sysmon ante inyección de procesos (\texttt{T1055}) y evasión de defensas. \\
\hline

% --- DE.AE (3 subcategorías) ---
\textbf{DE.AE}\newline Análisis de Eventos Adversos &
\textbf{\texttt{DE.AE-02}}: Los eventos potencialmente adversos se analizan mediante herramientas SIEM. &
\textbf{Tarea 1, Tarea 3 y Actividad 7} &
En la \textbf{Tarea 3} (conferencia \textit{Obfuscated Tale}) se estudió la diferencia entre sólo guardar logs y construir reglas analíticas de detección dentro de un \textbf{SIEM}. \\
\cline{2-4}\nopagebreak
 &
\textbf{\texttt{DE.AE-03}}: Se correlaciona información de múltiples fuentes. &
\textbf{Tarea 3 y Actividad 7} &
En la \textbf{Tarea 3} (conferencia \textit{Multi-Tenant Kill Chain Correlation}) se estudió la correlación de eventos entre identidades, sesiones SaaS y endpoints; en la \textbf{Actividad 7} se correlacionaron alertas de firewall, DLP y repositorios. \\
\cline{2-4}\nopagebreak
 &
\textbf{\texttt{DE.AE-08}}: Se declaran incidentes según criterios definidos considerando los falsos positivos. &
\textbf{Tarea 1, Tarea 3 y Actividad 7} &
En la \textbf{Tarea 1} se analizó cómo reducir el ruido y los \textbf{falsos positivos en el SOC} mediante alertas de alta fidelidad generadas por señuelos (\textit{Deception}), y en la \textbf{Tarea 3} se distinguió entre anomalía y evidencia de compromiso. \\
\hline
\end{longtable}
\endgroup

% ---------------------------------------------------------
\subsection{Función 5: RESPOND (RS) --- Responder (5 Subcategorías)}
% ---------------------------------------------------------

\begingroup
\scriptsize
\setlength{\tabcolsep}{4pt}
\renewcommand{\arraystretch}{1.24}
\begin{longtable}{|L{2.8cm}|L{4.5cm}|L{2.9cm}|L{6.5cm}|}
\caption{Mapeo de Categorías y Subcategorías de la función \textbf{RESPOND (RS)} con las actividades del semestre.}
\label{tab:mapeo-rs} \\
\hline
\rowcolor{csfRS}
\textbf{Categoría} & \textbf{Subcategoría Relacionada} & \textbf{Actividad / Práctica / Tarea} & \textbf{Qué se hizo en el semestre y cómo se relaciona} \\
\hline
\endfirsthead
\hline
\rowcolor{csfRS}
\textbf{Categoría} & \textbf{Subcategoría Relacionada} & \textbf{Actividad / Práctica / Tarea} & \textbf{Qué se hizo en el semestre y cómo se relaciona} \\
\hline
\endhead
\hline
\endfoot

% --- RS.MA (1 subcategoría) ---
\textbf{RS.MA}\newline Gestión de Incidentes &
\textbf{\texttt{RS.MA-03}}: Los incidentes se categorizan y priorizan (balanceando la recuperación rápida con observar al atacante en un Honeypot). &
\textbf{Tarea 1, Actividad 5 y Actividad 6} &
En la \textbf{Tarea 1} se estudió la estrategia de desviar y observar al atacante dentro de un \textbf{Honeypot} controlado para conocer sus herramientas antes de bloquearlo, y en las \textbf{Actividades 5 y 6} se categorizaron los incidentes por tipo. \\
\hline

% --- RS.AN (3 subcategorías) ---
\textbf{RS.AN}\newline Análisis de Incidentes &
\textbf{\texttt{RS.AN-03}}: Se realiza análisis para determinar qué ocurrió durante un incidente, su causa raíz y revisión de tecnología \textit{Deception}. &
\textbf{Tarea 1, Tarea 2 y Actividad 5} &
En la \textbf{Tarea 2} y en la \textbf{Actividad 5} se realizó el análisis de causa raíz y vector de ataque de incidentes reales (\textit{Operation Triangulation}, \textit{MOVEit}, \textit{EternalBlue}/\textit{WannaCry} y 5 ciberataques industriales). \\
\cline{2-4}\nopagebreak
 &
\textbf{\texttt{RS.AN-07}}: Los datos y metadatos del incidente se recolectan preservando su integridad. &
\textbf{Práctica 5 y Práctica 7} &
En la \textbf{Práctica 5}, al recolectar los archivos de contraseñas del servidor, se calculó su hash con \texttt{sha256sum} antes y después de la transferencia por red para preservar la integridad de la evidencia; en la \textbf{Práctica 7} se verificó el hash \texttt{SHA256} con \texttt{certutil}. \\
\cline{2-4}\nopagebreak
 &
\textbf{\texttt{RS.AN-08}}: Se estima y valida la magnitud de un incidente. &
\textbf{Práctica 4 y Práctica 6} &
En las \textbf{Prácticas 4 y 6} se ejecutaron herramientas de enumeración interna (\texttt{LinPEAS} y \texttt{escaner.sh}) para dimensionar el alcance del compromiso y descubrir otros equipos alcanzables en la red interna. \\
\hline

% --- RS.MI (1 subcategoría) ---
\textbf{RS.MI}\newline Mitigación de Incidentes &
\textbf{\texttt{RS.MI-02}}: Los incidentes son erradicados. &
\textbf{Práctica 4 y Actividad 7} &
En la \textbf{Práctica 4} se realizó la limpieza y eliminación de binarios compilados (\texttt{make clean}) tras la prueba de explotación, y en la \textbf{Actividad 7} se definió la revocación y rotación inmediata de llaves criptográficas comprometidas. \\
\hline
\end{longtable}
\endgroup

% ---------------------------------------------------------
\subsection{Función 6: RECOVER (RC) --- Recuperar (3 Subcategorías)}
% ---------------------------------------------------------

\begingroup
\scriptsize
\setlength{\tabcolsep}{4pt}
\renewcommand{\arraystretch}{1.24}
\begin{longtable}{|L{2.8cm}|L{4.5cm}|L{2.9cm}|L{6.5cm}|}
\caption{Mapeo de Categorías y Subcategorías de la función \textbf{RECOVER (RC)} con las actividades del semestre.}
\label{tab:mapeo-rc} \\
\hline
\rowcolor{csfRC}
\textbf{Categoría} & \textbf{Subcategoría Relacionada} & \textbf{Actividad / Práctica / Tarea} & \textbf{Qué se hizo en el semestre y cómo se relaciona} \\
\hline
\endfirsthead
\hline
\rowcolor{csfRC}
\textbf{Categoría} & \textbf{Subcategoría Relacionada} & \textbf{Actividad / Práctica / Tarea} & \textbf{Qué se hizo en el semestre y cómo se relaciona} \\
\hline
\endhead
\hline
\endfoot

% --- RC.RP (2 subcategorías) ---
\textbf{RC.RP}\newline Ejecución del Plan de Recuperación &
\textbf{\texttt{RC.RP-01}}: Se ejecuta la parte de recuperación del plan de respuesta a incidentes. &
\textbf{Actividad 5, Actividad 6 y Tarea 2} &
En la \textbf{Tarea 2} (\textit{WannaCry}) y en las \textbf{Actividades 5 y 6} se analizó el proceso de restauración operativa a partir de respaldos tras ataques de ransomware. \\
\cline{2-4}\nopagebreak
 &
\textbf{\texttt{RC.RP-03}}: La integridad de las copias de seguridad (respaldos) y activos de restauración se verifica antes de usarlos. &
\textbf{Práctica 5, Práctica 6 y Práctica 7} &
En la \textbf{Práctica 6} se inspeccionaron archivos de respaldo comprimidos (\texttt{tar -ztvf}), y en las \textbf{Prácticas 5 y 7} se aplicó la verificación criptográfica de integridad mediante hashes \textbf{SHA-256} (\texttt{sha256sum} y \texttt{certutil -hashfile}). \\
\hline

% --- RC.CO (1 subcategoría) ---
\textbf{RC.CO}\newline Comunicación de Recuperación &
\textbf{\texttt{RC.CO-03}}: El progreso en la restauración de capacidades operativas se comunica a las partes interesadas. &
\textbf{Actividad 5 y Actividad 7} &
En la \textbf{Actividad 5} se analizaron los comunicados emitidos por organizaciones afectadas tras restaurar sus servicios y mitigar brechas de seguridad. \\
\hline
\end{longtable}
\endgroup

% =========================================================
\section{Total de Categorías y Subcategorías Cubiertas en el Curso}
% =========================================================

Como resultado final de la revisión y mapeo de \textbf{The NIST Cybersecurity Framework (CSF) 2.0} frente a los temas vistos en clase, las Actividades 1 a 7, las Prácticas 1 a 7, la Presentación en clase y las Tareas 1 a 3, se encontró la siguiente cobertura:

\begin{itemize}[leftmargin=1.2cm, itemsep=3pt]
    \item \textbf{Categorías cubiertas en el curso:} Se encontraron \textbf{18 de las 22 categorías (81.8\%)} del framework relacionadas directamente con las actividades, prácticas y tareas realizadas en el semestre.
    \item \textbf{Subcategorías cubiertas en el curso:} Se encontraron \textbf{48 de las 106 subcategorías (45.3\%)} del framework cubiertas a lo largo del semestre.
\end{itemize}

\begin{table}[H]
\centering
\footnotesize
\setlength{\tabcolsep}{6pt}
\renewcommand{\arraystretch}{1.28}
\caption{Conteo final de Categorías y Subcategorías del NIST CSF 2.0 cubiertas en el curso.}
\label{tab:conteo-final}
\begin{tabular}{|L{4.2cm}|C{3.4cm}|C{3.4cm}|C{3.6cm}|}
\hline
\rowcolor{blue!20}
\textbf{Función Core del Framework} & \textbf{Categorías Cubiertas} & \textbf{Subcategorías Cubiertas} & \textbf{Porcentaje de Subcategorías} \\
\hline
\rowcolor{csfGV}
\textbf{1. GOVERN (\texttt{GV})} & \textbf{4 de 6} (66.7\%) & \textbf{11 de 31} & 35.5\% \\
\hline
\rowcolor{csfID}
\textbf{2. IDENTIFY (\texttt{ID})} & \textbf{3 de 3} (100.0\%) & \textbf{12 de 21} & 57.1\% \\
\hline
\rowcolor{csfPR}
\textbf{3. PROTECT (\texttt{PR})} & \textbf{4 de 5} (80.0\%) & \textbf{11 de 22} & 50.0\% \\
\hline
\rowcolor{csfDE}
\textbf{4. DETECT (\texttt{DE})} & \textbf{2 de 2} (100.0\%) & \textbf{6 de 11} & 54.5\% \\
\hline
\rowcolor{csfRS}
\textbf{5. RESPOND (\texttt{RS})} & \textbf{3 de 4} (75.0\%) & \textbf{5 de 13} & 38.5\% \\
\hline
\rowcolor{csfRC}
\textbf{6. RECOVER (\texttt{RC})} & \textbf{2 de 2} (100.0\%) & \textbf{3 de 8} & 37.5\% \\
\hline
\rowcolor{xlSummary}
\textbf{TOTAL EN EL CURSO} & \textbf{18 de 22 (81.8\%)} & \textbf{48 de 106} & \textbf{45.3\%} \\
\hline
\end{tabular}
\end{table}

% =========================================================
\newpage
\printbibliography[title={Referencias Bibliográficas}]
% =========================================================

\end{document}
